Na camada de segurança, os resultados do Capítulo~\ref{cap:desenvolvimento} respondem à questão colocada no Capítulo~\ref{cap:conceitos}, se a criptografia leve em \textit{software} compensa frente ao acelerador de AES presente nas placas. A implementação do Ascon-128a passou nos vetores oficiais, rejeitou as 4.907 adulterações aplicadas ao texto cifrado, aos dados associados, à \textit{tag} e ao \textit{nonce} nas duas placas, não realizou nenhuma alocação dinâmica, usou no máximo 436 bytes de pilha e cerca de 6,8~kB de \textit{Flash}, e manteve tempo de execução idêntico para cinco classes de entrada, o que corrobora a propriedade de tempo constante. A mensagem típica de 23 bytes é cifrada em 27~$\mu$s no ESP32 clássico e em 25~$\mu$s no ESP32-S3. Frente ao AES-128-GCM do próprio \textit{core}, o Ascon-128a foi mais rápido em todas as 1.200 execuções de cada placa, de 1,3 a 1,6 vez no ESP32 clássico e de 2,3 a 3,5 vezes no ESP32-S3, que é a placa de implantação. Nos dois chips a autenticação do GCM permanece em \textit{software} e o acelerador cobre só o bloco AES, e no ESP32-S3 o caminho por DMA tem custo fixo alto por chamada, o que explica a vantagem maior justamente onde ela importa.

A biblioteca \textit{AsconSecure} leva esse resultado ao enlace. Cada mensagem acrescenta 20 bytes, um contador de 4 bytes em claro e a \textit{tag} de 16, a chave nasce no próprio nó e nunca sai dele por outro caminho que não a bancada, e o receptor rejeita repetições. O custo de processamento da camada é desprezível diante do rádio, os 27~$\mu$s da cifragem contra cerca de 87~ms de tempo no ar na configuração adotada, de modo que o custo real de proteger o enlace está nos bytes adicionais transmitidos, e não no algoritmo escolhido. Esse custo é o mesmo para qualquer AEAD com \textit{tag} de 128 bits, e a escolha pelo Ascon-128a se sustenta na memória, no tamanho de código, no tempo constante e no tempo de execução, não em uma economia de energia de transmissão.
